\documentclass{stex}

\libusepackage{naproche}
\libinput{preamble}

\begin{document}
\begin{smodule}{omega-is-limit-ordinal.ftl}

\importmodule[libraries/set-theory]{omega.ftl}
\importmodule[libraries/set-theory]{limit-ordinals.ftl}

\begin{proposition*}[forthel]
  $\Infinity$ is a limit ordinal.
\end{proposition*}
\begin{proof}[forthel]
  $\Infinity$ is transitive.
  \begin{proof}
    Define $\Phi \DefEq \SClass{n}{\Infinity}{\text{ for all }m \Elem n\text{ we have }m \Elem \Infinity}$.

    (1) $\Zero \Elem \Phi$.

    (2) For all $n \Elem \Phi$ we have $\Succ(n) \Elem \Phi$.
    \begin{proof}
      Let $n \Elem \Phi$.
      Then every element of $n$ is contained in $\Infinity$.
      Hence every element of $\Succ(n)$ is contained in $\Infinity$.
      Thus $\Succ(n) \Elem \Phi$.
    \end{proof}

    Therefore $\Infinity \Subclass \Phi$.
    Consequently $\Infinity$ is transitive.
  \end{proof}

  Every element of $\Infinity$ is an ordinal.
  Hence every element of $\Infinity$ is transitive.
  Thus $\Infinity$ is an ordinal.

  $\Infinity$ is a limit ordinal.
  \begin{proof}
    Assume the contrary.
    We have $\Infinity \NotEq \Zero$.
    Hence $\Infinity$ is a successor ordinal.
    Take an ordinal $\alpha$ such that $\Succ(\alpha) \Eq \Infinity$.
    Then $\alpha \Elem \Infinity$.
    Thus $\Infinity \Eq \Succ(\alpha) \Elem \Infinity$.
    Contradiction.
  \end{proof}
\end{proof}

\end{smodule}
\end{document}
